% ============ 02. Contexto y guía de lectura ============
\section{Contexto y guía de lectura}

\subsection{Qué es este documento}

Este documento desarrolla el \textbf{criterio 2 de la rúbrica}:
\emph{Datos utilizados: suficiencia y corrección}. El repositorio auditado es el
Observatorio de investigadores reconocidos del Ministerio de Ciencia,
Tecnología e Innovación de Colombia (MinCiencias), que consolida y analiza
\textbf{6 convocatorias (2013--2021)} del padrón oficial de investigadores
y su producción científica. Esta serie cubre los \textbf{6 criterios} de la
rúbrica en documentos separados, más el documento maestro que los integra.

\textbf{En resumen:} Este criterio evalúa si los datos del proyecto eran suficientes y correctos. Conclusión principal: los datos correctos existen y son suficientes (30.086 investigadores únicos, 3.166.629 productos), pero el archivo consolidado versionado en el repositorio está incompleto: solo contiene 3 de las 6 convocatorias, por lo que un clon del repositorio no reproduce todos los resultados del informe.




\begin{tcolorbox}[nota]
\textbf{Cómo leer este documento:} cada PDF de esta serie desarrolla \emph{un}
criterio de la rúbrica. Los demás criterios se tratan a fondo en sus propios
documentos y en el \textbf{documento maestro} (\texttt{maestro\_auditoria.pdf}).
Si este es tu primer acercamiento, lee primero la portada y este contexto; al
final encontrarás las conclusiones del criterio.
\end{tcolorbox}

\subsection{Glosario (solo los términos que usa este documento)}

Esta tabla incluye únicamente los términos técnicos que aparecen en este
documento, explicados en lenguaje sencillo:

\begingroup
\small
\setlength{\tabcolsep}{3pt}
\begin{longtable}{>{\raggedright\arraybackslash}p{4.3cm}>{\raggedright\arraybackslash}p{9.8cm}}
\caption{Glosario de términos presentes en este documento}
\label{tab:glosario}\\
\toprule
\textbf{Término} & \textbf{Qué significa} \\
\midrule
\endfirsthead
\toprule
\textbf{Término} & \textbf{Qué significa} \\
\midrule
\endhead
    Matriz de transición & Tabla que muestra con qué probabilidad un investigador pasa de una categoría (p. ej. Junior) a otra (p. ej. Asociado) entre dos convocatorias. Permite ver si el sistema 'sube', 'mantiene' o 'pierde' investigadores.
 \\
    Mojibake & Texto ilegible por un problema de codificación (p. ej. 'Física' en lugar de 'Física'). Detectado en el CSV del repositorio.
 \\
    Contrato de esquema & Validación automática que garantiza que un archivo de datos tenga las columnas, tipos y valores esperados antes de analizarlo.
 \\
    Pipeline & Cadena de procesos que lleva los datos desde la fuente hasta el análisis (ingesta, limpieza, transformación, modelado). También llamada 'cadena de procesamiento' o 'flujo de datos'.
 \\
\bottomrule
\end{longtable}
\endgroup
